% ═══════════════════════════════════════════════════════════════════
% ЧАСТЬ XV. ДИСКРЕТНЫЙ ПОТОК, БАШНИ, ЗАКЛЮЧЕНИЕ
% ═══════════════════════════════════════════════════════════════════
\part{Дискретный поток, башни, заключение}

\section{Дискретный поток на шахматной сетке}\label{sec:flow}

Сертификат~E доказал завершимость потока в общей форме
(теорема~\ref{th:E}); здесь мы разбираем механику потока
подробно --- на уровне состояний, переходов и орбит.

\subsection{Состояния и переходы}

Состояние потока --- пара $(p,\varphi)$: позиция $p$ на решётке
$W\times H$ и фаза $\varphi\in\mu_4$ (одна из четырёх фаз поворота).
Переход $(p,\varphi)\to(p',\varphi')$ определяется шагом
$(a,b)$ и правилом поправки: позиция сдвигается по модулю сетки,
фаза умножается на генератор поворота. Пространство состояний
конечно: $4WH$ состояний --- первая посылка завершимости.

Моновариант открытия --- множество посещённых ячеек (без фазы):
каждый переход либо открывает новую ячейку (моновариант растёт),
либо попадает в посещённую --- и тогда по детерминированности
переходов поток входит в цикл. Ограниченность моноварианта
($WH$ ячеек) и строгий рост в режиме открытия дают завершимость за
конечное время --- вторая и третья посылки теоремы.

\subsection{Точное время и его проверка}

Время завершимости $t^*=\lcm(W/\gcd(a,W),\,H/\gcd(b,H))$ ---
момент первого совместного возврата позиции. Для эталонного
конвейера $48\times48$ с шагом $(1,1)$: $t^*=48$ --- и поток
действительно завершается за $48$ шагов, посетив ровно $48$
ячеек, с $212$ событиями трения, $432$ краевыми ячейками и цепным
кодом длины $114$ --- координатная сверка $48/212/432/114$
полная. Это редкий случай, когда теорема даёт не только
существование, но и точное значение, подтверждённое независимым
конвейером (раздел~\ref{sec:binary}).

\subsection{Терминальный цикл как $\mu_4$-орбита}

Терминальное состояние потока --- цикл; его структура вычисляется
явно: цикл длины $4$ --- в точности $\mu_4$-орбита шага. Фазовый
поворот переводит конфигурацию в конфигурацию, и после четырёх
переходов и позиция, и фаза возвращаются. Таким образом,
«циклическая структура» терминального состояния --- не эмпирика,
а следствие представления $\mu_4$: сертификат~E-4 фиксирует это
как PASS с прямым перебором орбит.

\subsection{Масштабирование и слои}

Формула $t^*$ проверена на масштабах $48\to96\to192\to384$ ---
время растёт линейно с размером сетки при единичном шаге, что
согласовано с формулой $\lcm$. Слои верхних стендов:
K3-слой (отражения с торможением) завершим на тех же правилах;
Клейн-слой добавляет Зингер-порождение порядка $7$ --- время
завершимости умножается на $7/\gcd(7,H)$, и при $H$ кратном $7$
поток замыкается за исходное время. Все восемь пунктов протокола
E --- PASS (таблица раздела~\ref{sec:protoE}).

\section{Теория башен больших матриц}\label{sec:towers}

\subsection{Блочно-кронекерова конструкция}

Башни строятся из базового блока $A_0$ размера $m_0$ повторным
кронекеровым возведением с фазовым сдвигом:
$A_{k+1}=A_k\otimes I + I\otimes A_k^{(\varphi)}$, где верхний
индекс обозначает поворот фазы. Для операторов $O_\chi$ базовый
блок $28\times28$; башня $O_\chi(28k)$ имеет размер $28k$ при
$k\in\{1,2,4,8,16\}$ --- уровни $28,56,112,224,448$.

\subsection{Спектральная стабильность}

Колебания собственных значений по башне --- максимум
$2.06\cdot10^{-16}$ на нижнем уровне и монотонное убывание до
$4.2\cdot10^{-17}$ --- показывают, что кронекерово расширение не
вносит спектрального дрейфа: собственные значения стабилизируются
на машинной точности. Причина --- блочная структура: спектр башни
есть объединение спектров блоков с фазовыми множителями из
конечной группы, и повторные уровни не добавляют новых величин, а
лишь повторяют существующие с большей кратностью.

\subsection{Разреженные Грам-башни}

Грам-башни $G_N$ растут как $22k$: уровни $22,220,2200,22000,
220000$. Формат CSR хранит только ненулевые элементы --- плотность
ненулевых падает как $O(1/N)$, поэтому память растёт линейно
(~39 МБ на уровне $220\,000$), а сборка остаётся субсекундной.
Краевой спектр $[-3.9890438,\,1.0000000]$ повторяется на всех
уровнях с пятью знаками --- спектральные края башни фиксированы
базовым блоком, и большие матрицы не приносят новых краевых
значений. Это ключевое свойство для практических применений:
свойства бесконечного предела читаются на малых уровнях.

\begin{databox}[практический вывод о башнях]
Спектральная стабильность + линейная память + субсекундная сборка
= башни масштабируются на кластерные платформы без изменения
формата. Блочный формат готов к GPU-переносу: каждый блок
независим, коммуникация минимальна.
\end{databox}

\section{Заключение: полная арка программы}\label{sec:final}

Монография свела программу в единую арку. Кратко пройдём её
ещё раз --- теперь сверху.

\subsection{Что построено}

\emph{Основание} --- динамический принцип: геометрия поставляет
только целочисленную тройку $(n,B,\lambda_0)$; вся структура
строится алгеброй. Принцип превращает фреймворк в конвейер без
свободных параметров: спиновая фаза $\pi/N$, торможение
$\delta^4/k$, эффективная фаза $\delta^5/k$, объединённая формула
$\Delta_{Ch}$ --- всё выведено и всё проверяемо.

\emph{Калибровка} --- три стенда: тор с полным явным выводом
тройки $(4,1,4\pi^2)$ и точной $\mu_4$-эквивариантностью;
K3 с максимальным рангом $\rho=20$, сигнатурой $(1,19)$ и
определителем $64$ --- всё над $\QQ$ по явной матрице $49\times49$;
Клейн с арифметикой гептады $\Delta=-7^3$, $j=-3375$ и
восстановлением $j$ из периода с $57$ знаками.

\emph{Сертификаты} --- восемь документов от A до H: явный дивизор,
закрытие ядра $29$ периодами, эквивариантность $22=1+7+7+7$,
универсальная $\mu_4$-теорема, завершимость потока с точным
временем, рационализация с априорной границей $Q=256$, индексы
решётки Ходжа со слоем Коула--Селберга, и --- вершина --- стенд
$N=15/30$ в полной $\Ga$-нормировке Гросса: SNF-тип
$(1,1,5,5,15,15,15,15)$, дискриминант $1125^2$, объём
$\mathrm{vol}_h=1125$, граница знаменателей $480$.

\emph{Замкнутая система} --- стенд $N=15/30$: восемь лемм и шесть
теорем с полными доказательствами, от комбинаторики циклов до
поляризации и индексов; каждая величина --- в замкнутой форме,
каждое тождество --- либо доказано, либо сертифицировано
независимым интегрированием до $10^{-36}$; протокол V1--V9
воспроизводит всё одной командой.

\subsection{Чем программа отличается}

Три черты определяют отличие программы от стандартной
вычислительной практики. Во-первых, \emph{замкнутость}: ни одна
величина не берётся извне без вывода или независимого сертификата;
даже классические теоремы (Шиода, Рибет), используемые как рамки,
сопровождаются точными целочисленными проверками внутри текста.
Во-вторых, \emph{воспроизводимость}: полный протокол, пять языков,
ядро Lean --- согласование независимых реализаций до целых
значений и до машинных пределов. В-третьих, \emph{переносимость}:
новая геометрия входит тройкой, и конвейер работает без изменений
--- лестница от тора до $N=30$ есть доказательство переносимости
конструкцией.

\subsection{Куда идёт программа}

Лестница продолжается: уровни $\pi/N$ с ростом $N$, блоки
проводников с ростом числа индексов, башни матриц с ростом
масштаба. Инфраструктура --- лаборатория, сертификаты, Lean-панель,
мультиязычная верификация --- спроектирована так, чтобы каждая
новая ступень вставала на готовое основание: конструктор
экспериментов лаборатории позволяет запускать новые конфигурации
без изменения кода, а протокол --- сертифицировать их теми же
правилами. Именно в этом смысле монография --- не итог, а
фундамент: арка построена, и каждая новая ступень получает своё
место в уже работающей системе.
